\documentclass[12pt,a4paper]{article}
\usepackage{amsmath,amssymb}

\title{Stability of a Feedback Loop}
\author{Viktor Hrabal}
\date{May 12, 2016}

\newcommand{\Lap}{\mathcal{L}}
\newcommand{\sv}{s}

\begin{document}
\maketitle

\section{A first-order plant}
Consider a plant with transfer function $G(\sv)=\dfrac{K}{\tau\sv+1}$, where $K$ is the DC gain and $\tau>0$ the time constant. Closing the loop with unity feedback gives
\begin{equation}
  T(\sv)=\frac{G(\sv)}{1+G(\sv)}=\frac{K}{\tau\sv+1+K}.
\end{equation}
The single pole is at $\sv=-(1+K)/\tau$, which lies in the left half-plane for all $K>-1$.

\section{Adding delay}
Real systems have delay. Replacing $G$ by $G(\sv)e^{-\sv d}$ with delay $d>0$ makes the characteristic equation transcendental:
\[
  1+\frac{K\,e^{-\sv d}}{\tau\sv+1}=0 .
\]
On the imaginary axis $\sv=\mathrm{i}\omega$ the magnitude condition reads
\[
  \frac{K}{\sqrt{1+\tau^2\omega^2}}=1,
\]
which has the solution $\omega_c=\sqrt{K^2-1}/\tau$ for $K>1$.

\section*{Procedure for a quick check}
\begin{quote}
Plot the open-loop frequency response and find the frequency at which the gain is one. The phase margin is $180^\circ$ plus the phase at that frequency.
\end{quote}

The phase of $G(\mathrm{i}\omega)e^{-\mathrm{i}\omega d}$ at $\omega_c$ equals $-\arctan(\tau\omega_c)-\omega_c d$. The loop is stable if this phase is above $-\pi$, that is, if
\[
  d < \frac{\pi-\arctan(\tau\omega_c)}{\omega_c}.
\]

\begin{flushleft}
Example values:
\end{flushleft}
\begin{center}
\begin{tabular}{c|c|c}
$K$ & $\tau$ (s) & $d_{\max}$ (s)\\
\hline
2 & 1 & 1.37\\
5 & 1 & 0.68\\
5 & 0.5 & 0.34\\
\end{tabular}
\end{center}

Larger gain thus reduces the tolerable delay, a trade-off familiar to every control engineer.

\end{document}
